\documentclass[11pt]{article}
\usepackage[utf8]{inputenc}
\usepackage[T1]{fontenc}
\usepackage[english]{babel}
\usepackage{amsmath,amssymb,amsthm}
\usepackage{booktabs}
\usepackage{geometry}
\usepackage{hyperref}
\usepackage{longtable}
\usepackage{natbib}
\usepackage{url}
\geometry{margin=1in}

\title{The IDEAV Quintet: A Metadata-Compiled Data Model for Low-Code Business Applications}
\author{Aleksej P. Semenov\\Integram / IDEAV}
\date{Draft for arXiv cs.DB, May 2026}

\newtheorem{definition}{Definition}
\newtheorem{theorem}{Theorem}

\begin{document}
\maketitle

\begin{abstract}
Entity--Attribute--Value (EAV) schemas are often treated as a necessary but
dangerous compromise: they make application schemas extensible, yet they move
typing, integrity, navigation, and query planning out of the relational model
and into ad hoc application code. This paper describes the IDEAV quintet, a
metadata-compiled data model used in the Integram low-code platform to retain
EAV-like schema plasticity while making the missing semantics explicit. The
model represents all business metadata and values inside a uniform tuple space,
but separates their roles into five cooperating elements: types, requisites,
objects, links, and a compiler. The compiler, implemented by the
\texttt{Compile\_Report} routine, turns a user-level report definition into SQL
by following four formally defined navigation cases: reference to the current
object, child array of the current object, reference from the new object back to
the current object, and parent array reached from the current object. We give a
formal description of the model, a proof sketch that the compiler covers every
connected report graph expressible by the metadata, and a storage discussion
grounded in the single-table patent family RU2650032C1 and US11138174B2. The
paper also compares IDEAV with medical EAV systems, business EAV failures, and
modern low-code products. Because controlled public benchmarks are not yet
available, performance statements are stated as hypotheses and as a benchmark
protocol rather than as measured claims.
\end{abstract}

\noindent\textbf{Keywords:} low-code, metadata-driven database, EAV, schema evolution, query compilation.

\section{Introduction}

Low-code platforms promise that a business analyst can change a process without
waiting for a database migration. Classical relational design promises the
opposite virtue: data structures are declared explicitly, typed, normalized, and
optimized before workloads depend on them \citep{codd1970relational}. The
practical tension between these promises is the reason Entity--Attribute--Value
(EAV) schemas keep reappearing. EAV is easy to extend, but unstructured EAV
frequently degenerates into dynamic string storage where the database can no
longer see the real schema \citep{karwin2010sql,kyte2003effective}.

The usual question is therefore: why not just use EAV? The short answer is that
EAV by itself is not a model. It is a physical encoding pattern. It says that
attributes may be rows, but it does not say how attributes are typed, how
relations are distinguished from scalar values, how a query compiler should
reconstruct joins, how arrays are represented, or what application invariant
prevents a collection of triples from becoming a private programming language.

IDEAV starts from the same need for late-bound structure, but it does not stop at
attribute rows. Its central contribution is a quintet:

\[
  \mathcal{Q}_{IDEAV} =
  \langle \mathcal{T}, \mathcal{R}, \mathcal{O}, \mathcal{L}, \mathcal{C} \rangle ,
\]

where \(\mathcal{T}\) is the set of logical types, \(\mathcal{R}\) is the set of
requisites (fields and relations), \(\mathcal{O}\) is the set of objects,
\(\mathcal{L}\) is the set of typed links, and \(\mathcal{C}\) is a metadata
compiler. This fifth element is important. Without the compiler, the first four
elements still risk becoming a named variant of EAV. With the compiler, reports,
forms, and API calls operate against a graph of declared metadata rather than
against arbitrary key-value strings.

The model has been developed in the commercial Integram platform since 2015,
with conceptual roots dating to 2003. It is also related to the patent family
``Electronic database and method of its formation'' \citep{semenov2018ru} and
``Electronic database and method for forming same'' \citep{semenov2021us}. The
patents are not used here as marketing evidence. They are cited to establish a
documented design lineage for a single-table, role-coded representation of
metadata and data.

\section{Related Work}

\subsection{Medical and Scientific EAV}

The strongest academic tradition around EAV is not commerce but medicine and
scientific data capture. ACT/DB used attribute-value storage to manage clinical
trial variables whose structure changed over time \citep{nadkarni1998actdb}.
EAV/CR extended the idea with classes and relationships for heterogeneous
scientific data \citep{nadkarni1999eavcr,marenco2003eavcr}. Web clinical
databases and i2b2 showed that ontology-like metadata can support rich,
evolving datasets \citep{anhoj2003generic,murphy2010i2b2}.

These systems answer part of the ``why not EAV'' question: EAV can work when the
domain accepts metadata as a first-class artifact. They also reveal the missing
piece for operational business software. A hospital research schema often
prioritizes sparse observation capture and later analysis. A CRM or ERP system
also needs daily transactional forms, permissions, references, child tables,
bulk imports, and user-defined reports. IDEAV therefore treats metadata not only
as descriptive schema, but as executable schema for a compiler.

\subsection{Business Uses and Failure Modes}

Business systems adopted EAV for product attributes, custom fields, and CRM-like
extension points. The cautionary literature is familiar: ``Bad CaRMa'' describes
the pain of a catalog system whose entity-attribute design hid useful structure
from the database engine \citep{gorman2006badcarma}. SQL antipattern discussions
make similar points: generic attribute rows often defeat constraints, make
queries verbose, and encourage application-only integrity \citep{karwin2010sql}.

Magento and WooCommerce illustrate a more nuanced pattern. Adobe Commerce still
documents extensible attributes \citep{adobeCommerceAttributes}, while
WooCommerce introduced product lookup tables and high-performance order storage
to reduce dependence on broad metadata tables for hot paths
\citep{woocommerceLookup,woocommerceHPOS}. These are not simply failures of EAV;
they are evidence that an extensibility encoding without a formal compiler,
typed navigation rules, and workload-aware projections tends to require
secondary structures later.

The present draft does not cite unsupported anecdotal cases, including
unarchived stories sometimes retold about cascade deletion in CRM products, as
primary evidence. That omission is deliberate. The argument is stronger when it
rests on public documentation and reproducible model properties.
Likewise, the requested Kustomer comparison is retained only as a candidate for
future archival verification: without a primary public source, it should not
carry an empirical claim.

\subsection{Modern Low-Code Platforms}

Low-code systems solve schema evolution with different trade-offs. Some expose a
spreadsheet-like database; some connect to an external relational database; some
generate conventional tables; some keep metadata inside an application runtime.
Table~\ref{tab:lowcode} summarizes the comparison at the model level.

\begin{longtable}{p{0.18\linewidth}p{0.31\linewidth}p{0.39\linewidth}}
\caption{Selected low-code systems and their schema strategy.}\label{tab:lowcode}\\
\toprule
System & Storage strategy & Implication for the EAV question \\
\midrule
\endfirsthead
\toprule
System & Storage strategy & Implication for the EAV question \\
\midrule
\endhead
Airtable & Hosted base/table abstraction with API access \citep{airtableApi} &
Strong user experience, but storage and compilation semantics are proprietary. \\
Bubble & Application data types managed by the Bubble runtime
\citep{bubbleData} & Flexible, but the database model is not a portable formal
contract. \\
Retool & Connects forms and workflows to external resources
\citep{retoolResources} & Keeps real databases available; schema evolution is
delegated to connected systems. \\
NocoDB, Baserow & Spreadsheet-like interfaces over structured backends
\citep{nocodb,baserow} & Useful low-code layer, usually table-oriented rather
than a single formal metadata tuple space. \\
Directus, Strapi & Headless data/content platforms over database schemas
\citep{directus,strapi} & Excellent when the schema exists; less focused on
compiling arbitrary business metadata into joins. \\
Caspio, Knack, AppSheet & Hosted app builders \citep{caspio,knack,appsheet} &
Reduce development cost, but do not publish a formal equivalent of the IDEAV
compiler. \\
Integram / IDEAV & Single role-coded tuple space plus metadata compiler
\citep{integramWorkflow2026,integramMcp2026} & Treats schema evolution and query
navigation as one formal problem. \\
\bottomrule
\end{longtable}

\subsection{Database Systems Context}

IDEAV is not a replacement for relational algebra, NewSQL, or document stores.
It is a higher-level metadata model that uses a relational engine as execution
substrate. NewSQL work shows that relational systems can scale transactional
workloads when the schema is known \citep{pavlo2016newsql}. NoSQL and polyglot
persistence work shows that alternate models can simplify specific access
patterns \citep{sadalage2012nosql,floratou2012nosql}. IDEAV asks a narrower
question: if business users create schema at runtime, what formal metadata must
exist so that the generated relational query is still predictable?

\section{The IDEAV Quintet Model}

\subsection{Physical Tuple Space}

The patented core describes an electronic database where elements are organized
in a single table with identifiers, parent identifiers, values, and a type code
\citep{semenov2018ru,semenov2021us}. The Integram implementation augments this
minimal tuple with metadata columns used by the compiler. We write the physical
universe as:

\[
  U = \{q_i = (id_i, t_i, up_i, val_i, A_i)\},
\]

where \(id_i\) is a stable identifier, \(t_i\) is the type identifier, \(up_i\)
is the parent/object identifier, \(val_i\) is the stored value, and \(A_i\) is a
finite map of implementation attributes such as aliases, reference markers,
array markers, uniqueness flags, and permissions. This is not ``just EAV''
because \(t\) and \(up\) are interpreted by declared role metadata rather than
by each application screen independently.

\subsection{Five Roles}

\begin{definition}[Type]
A type \(\tau \in \mathcal{T}\) declares a logical business class. In the API it
is created by \texttt{\_d\_new}. Its first name field can be addressed as
\texttt{t\(\tau\)}.
\[
  \tau = (id, name, parent, unique, permissions).
\]
\end{definition}

\begin{definition}[Requisite]
A requisite \(r \in \mathcal{R}\) is a field or structural relation belonging to
a type. It is created by \texttt{\_d\_req}. A scalar requisite stores values; a
reference requisite contains \texttt{ref} or \texttt{ref\_id}; an array
requisite contains \texttt{arr\_id}.
\[
  r = (id, owner, name, domain, ref, ref\_id, arr\_id, attrs).
\]
\end{definition}

\begin{definition}[Object]
An object \(o \in \mathcal{O}\) is a typed business record.
\[
  o = (id, type(o), parent(o), name(o), \{r \mapsto v_r\}).
\]
Scalar values are stored as tuples under the object; the record name is stored
as the value of the type row itself.
\end{definition}

\begin{definition}[Link]
A link \(\ell \in \mathcal{L}\) is a typed navigation fact between objects or
between a parent object and a child array item. References and arrays are kept
separate because they compile into different join predicates.
\[
  \ell_{ref}: o_s \xrightarrow{r} o_t,\qquad
  \ell_{arr}: o_p \xrightarrow{a} o_c.
\]
\end{definition}

\begin{definition}[Compiler]
The compiler \(\mathcal{C}\) maps metadata and a report definition to SQL:
\[
  \mathcal{C}: (\mathcal{T},\mathcal{R},\mathcal{L}, Report, \Gamma)
  \rightarrow SQL ,
\]
where \(\Gamma\) contains aliases, filters, functions, and current block
context. In the implementation this role is represented by
\texttt{Compile\_Report}.
\end{definition}

\subsection{The Metadata Compiler}

A report is a list of requested columns. Each column belongs to a logical type.
The compiler maintains the set \(J\) of logical types already connected to the
SQL query. For every new requested type \(n\), it searches for a type
\(j \in J\) that can connect to \(n\). The implementation documents four cases:

\begin{center}
\begin{tabular}{lll}
\toprule
Case & Metadata condition & Join shape \\
\midrule
Reference to us & \(j\) has reference requisite to \(n\) &
\(j.ref = n.id\) \\
We are an Array & \(n\) is an array child of \(j\) &
\(n.up = j.id\) \\
We have a Reference & \(n\) has reference requisite to \(j\) &
\(n.ref = j.id\) \\
We got an Array & \(j\) is an array child of \(n\) &
\(j.up = n.id\) \\
\bottomrule
\end{tabular}
\end{center}

In simplified form:

\begin{verbatim}
J := {master type from the first report column}
while uncompiled columns remain:
    progress := false
    for column c with type n not in J:
        for each j in J:
            if j references n: add join "j.ref = n.id"
            else if n is array child of j: add join "n.up = j.id"
            else if n references j: add join "n.ref = j.id"
            else if j is array child of n: add join "j.up = n.id"
            if a case matched:
                J := J union {n}
                progress := true
                break
    if not progress: fail with "cannot connect report columns"
\end{verbatim}

This is where IDEAV diverges from ordinary EAV. A generic attribute table can
store ``customer'', ``invoice'', and ``invoice lines'' as rows, but it does not
necessarily know which join is legal. IDEAV makes the four navigations explicit
metadata and requires the compiler to choose one.

\subsection{Correctness Sketch}

\begin{theorem}[Connected report compilation]
Let \(G_M=(V,E)\) be the undirected graph induced by IDEAV reference and array
metadata for a report, where vertices are logical types and edges are legal
instances of one of the four compiler cases. If the requested report types form
a connected subgraph of \(G_M\) and every edge has an unambiguous metadata role,
then \(\mathcal{C}\) can add a join path for every requested type.
\end{theorem}

\begin{proof}
Choose the first report type as the master vertex. The invariant after each
successful compiler step is that \(J\) is a connected set of vertices and the
generated SQL contains a join path from the master to every vertex in \(J\).
Initially the invariant holds for the singleton master. Suppose it holds for
\(k\) vertices. Because the requested subgraph is connected, if not all vertices
are in \(J\), there exists an edge from some \(j \in J\) to some
\(n \notin J\). By construction every metadata edge is exactly one of the four
compiler cases, so the compiler can add the corresponding join predicate and
extend \(J\) to \(J \cup \{n\}\). The invariant is preserved. By induction over
the number of requested types, all vertices are joined. If no such edge exists,
the report is disconnected or ambiguous, and failure is the correct behavior.
\end{proof}

The theorem does not claim cost optimality. It claims semantic coverage for the
metadata graph that the platform itself allows users to create. Query cost is an
execution-engine and workload question, not a proof consequence.

\subsection{Physical Storage and Indexing}

The storage design keeps metadata and data in one physical relation, commonly
with composite indexes over identifiers, type, parent, and value
\citep{integramWorkflow2026}. This fixed tuple shape has two engineering
advantages: schema changes become data changes, and the compiler can target a
stable set of indexed columns. However, it would be misleading to treat this
alone as a fixed performance bound. B-tree height depends on row count, page
size, key distribution, and partitioning strategy. The defensible claim is
more specific: a fixed tuple schema and partitionable identifier ranges give the
implementation a stable mechanism for keeping metadata operations and generated
joins inside a known operational envelope. Controlled benchmarks are needed to
turn that mechanism into measured latency claims.

\section{Empirical Validation}

The model should be evaluated at three levels.

\paragraph{Compiler correctness.}
Synthetic metadata graphs should cover every legal combination of references and
arrays. For each graph, a generated report should be compared with an equivalent
hand-written relational query over a normalized schema. The minimum test suite
contains chains, stars, diamonds, reverse references, parent-to-child arrays,
child-to-parent arrays, and disconnected reports that must fail.

\paragraph{Schema evolution.}
A workload should create new types, add scalar requisites, add references, add
arrays, and then execute reports without DDL migrations. The comparison systems
should include conventional migrations, a generic EAV table, a JSON/document
column design, and a metadata-compiled IDEAV design.

\paragraph{Performance.}
Benchmarks should report write latency, report compilation time, generated SQL
execution time, index size, and storage overhead as functions of:

\[
  N_o = \text{objects},\quad
  N_t = \text{types},\quad
  N_r = \text{requisites per type},\quad
  d = \text{report join depth}.
\]

Until those results are public, production history since 2015 should be treated
as evidence of feasibility, not as a controlled performance proof. The benchmark
protocol above is therefore part of the contribution: it identifies what must
be measured before IDEAV can be compared fairly with relational, EAV, and JSON
alternatives.

\section{Discussion}

IDEAV reframes the usual EAV debate. The problem with EAV in business systems is
not that attributes are rows. The problem is that the row encoding is commonly
asked to carry schema, constraints, navigation, UI behavior, and reporting
semantics without a formal layer. IDEAV assigns those responsibilities to
metadata and then makes the metadata executable through a compiler.

This design also changes the role of AI-assisted development. A language model
or scripted agent can call a compact API such as \texttt{\_d\_new},
\texttt{\_d\_req}, \texttt{\_d\_ref}, \texttt{\_m\_new}, and
\texttt{\_m\_save}. Those calls do not merely create tables in a UI; they extend
the same metadata graph that the compiler later uses for reports and forms. The
AI surface is therefore smaller than a full SQL migration surface but more
formal than a screen-scraping automation layer.

The limitations are real. Generated joins may be semantically correct and still
slow. Highly analytical workloads may need materialized projections. Strong
typing, migration history, and multi-tenant isolation must be specified
carefully. The model also depends on good metadata governance: if users create
poor types and relations, a compiler can preserve their meaning but cannot make
the domain design good.

\section{Conclusion}

EAV failed in many business settings because it was used as a shortcut around
modeling rather than as one layer in a formal model. The IDEAV quintet proposes
a different contract: keep the extensible tuple space, but require explicit
types, requisites, objects, links, and a compiler. The result is a low-code data
model in which schema evolution is cheap, but not invisible; relations are
dynamic, but not arbitrary; and generated SQL follows metadata rules that can be
stated, tested, and reviewed.

The next research step is empirical. Public benchmarks should compare IDEAV with
classical EAV, normalized relational schemas, and JSON/document designs on
schema evolution, report compilation, and query execution. If those measurements
confirm the operational hypotheses, IDEAV becomes not a rejection of relational
databases, but a disciplined metadata layer above them.

\appendix

\section{Patent References}

RU2650032C1, ``Electronic database and method of its formation'', lists Aleksej
Petrovich Semenov as inventor and describes a database whose elements are
organized in a single table containing identifiers, parent identifiers, values,
and datatype codes \citep{semenov2018ru}. The priority date is 2017-03-20 and
the publication date is 2018-04-06.

US11138174B2, ``Electronic database and method for forming same'', belongs to
the same patent family and describes the single-table organization with rows
for root element, datatype, term, term attribute, and data
\citep{semenov2021us}. The priority date is 2017-03-20 and the publication date
is 2021-10-05.

The patents establish novelty claims around the physical organization. This
paper's separate claim is model-level: IDEAV should be evaluated as the
combination of physical tuple space, formal metadata roles, and compiler.

\section{IDEAV API Sketch}

The public repository documentation exposes a compact API for creating metadata
and data \citep{integramMcp2026}. The following examples are schematic:

\begin{verbatim}
# Create a logical type/table.
POST _d_new
  val = "Customer"

# Add a scalar requisite/column to a type.
POST _d_req/{customerTypeId}
  val = "Phone"
  t   = "{stringTypeId}"

# Create a reference domain and attach it as a requisite.
POST _d_ref/{targetTypeId}
POST _d_req/{sourceTypeId}
  val = "Customer"
  t   = "{createdReferenceTypeId}"

# Attach an existing type as a subordinate array.
POST _d_req/{parentTypeId}
  val = "Invoice lines"
  t   = "{lineTypeId}"

# Create and save objects.
POST _m_new/{typeId}
POST _m_save/{objectId}
  t{typeId} = "Object name"
  {requisiteId} = "value"

# Inspect metadata.
GET metadata?JSON=1
\end{verbatim}

These commands are important because they demonstrate that IDEAV schema editing
is not a separate migration subsystem. The API mutates the same metadata graph
that the report compiler reads.

\bibliographystyle{plainnat}
\bibliography{references}

\end{document}
